\documentclass[11pt]{article}

\usepackage[margin=1in]{geometry}
\usepackage{enumitem}
\usepackage{amsmath, amssymb}

\setlist[enumerate]{itemsep=0.6em, topsep=0.6em}

\begin{document}

\begin{center}
  {\Large \textbf{PHYS 320 Examination}}\\[0.25em]
  {\large \textbf{Thermodynamics and Statistical Mechanics -- Answer Key}}\\[0.5em]
\end{center}

\begin{enumerate}[label=\textbf{\arabic*.}]

  \item \textbf{(C)} Gibbs free energy $G$.\\
  At constant temperature and pressure (most common lab conditions), the Gibbs free energy 
  $G = H - TS$ is minimized at equilibrium. The condition $dG = 0$ determines spontaneity: 
  $\Delta G < 0$ for spontaneous processes.

  \item \textbf{(B)} $TV^{\gamma-1} = \text{constant}$.\\
  For a reversible adiabatic process ($Q = 0$), the relationship $TV^{\gamma-1} = \text{constant}$ 
  holds, where $\gamma = C_P/C_V$ is the heat capacity ratio. Equivalently, $PV^\gamma = \text{constant}$.

  \item \textbf{(C)} A sum over Boltzmann factors for all states, encoding thermodynamic information.\\
  The partition function $Z = \sum_i e^{-\beta E_i}$ (where $\beta = 1/k_BT$) is the central object 
  in statistical mechanics. All thermodynamic quantities can be derived from $Z$, including free 
  energy $F = -k_B T \ln Z$.

  \item \textbf{(A)} $\eta = 1 - \frac{T_C}{T_H}$.\\
  The Carnot efficiency represents the maximum possible efficiency for any heat engine operating 
  between two reservoirs. It depends only on the temperatures and sets a fundamental limit based 
  on the second law of thermodynamics.

  \item \emph{Sample answer:}\\
  \textbf{Zeroth Law}: If two systems are each in thermal equilibrium with a third system, they are 
  in thermal equilibrium with each other. This establishes the concept of temperature as a transitive 
  property.

  \textbf{First Law}: $\Delta U = Q - W$ (energy conservation). The change in internal energy equals 
  heat absorbed minus work done. This states that energy cannot be created or destroyed, only converted.

  \textbf{Second Law}: In an isolated system, entropy never decreases: $\Delta S \geq 0$. This gives 
  directionality to time and explains why certain processes are irreversible. It limits the efficiency 
  of heat engines and prohibits perpetual motion machines of the second kind.

  \textbf{Third Law}: As temperature approaches absolute zero, the entropy of a perfect crystal approaches 
  zero: $S \to 0$ as $T \to 0$. This provides an absolute scale for entropy and implies that absolute 
  zero is unattainable.

  Together, these laws constrain energy conversion, establish equilibrium conditions, and explain 
  irreversibility in nature.

  \item \emph{Sample answer:}\\
  \textbf{Boltzmann's formula}: $S = k_B \ln \Omega$, where $\Omega$ is the number of accessible 
  microstates corresponding to a macrostate.

  \textbf{Derivation concept}:
  \begin{itemize}
    \item A macrostate with more microstates is more probable
    \item Entropy must be additive for independent systems: if system 1 has $\Omega_1$ states and 
    system 2 has $\Omega_2$ states, the combined system has $\Omega = \Omega_1 \times \Omega_2$ states
    \item Additivity requires $S(\Omega_1 \times \Omega_2) = S(\Omega_1) + S(\Omega_2)$, which is 
    satisfied by the logarithm: $\ln(\Omega_1 \times \Omega_2) = \ln \Omega_1 + \ln \Omega_2$
  \end{itemize}

  \textbf{Directionality}: Spontaneous processes evolve toward macrostates with larger $\Omega$ 
  (more microstates), increasing entropy. This explains why systems naturally progress from ordered 
  to disordered states and why time has a preferred direction (thermodynamic arrow of time).

  \item \emph{Sample answer:}\\
  \textbf{Isothermal process} (constant $T$):
  \begin{itemize}
    \item For ideal gas: $\Delta U = 0$ (internal energy depends only on $T$)
    \item $Q = W = nRT \ln(V_f/V_i)$ for expansion
    \item Entropy increases: $\Delta S = nR \ln(V_f/V_i) > 0$ for expansion
  \end{itemize}

  \textbf{Adiabatic process} ($Q = 0$):
  \begin{itemize}
    \item Temperature decreases during expansion: $T_f < T_i$
    \item $\Delta U = -W$ (work comes from internal energy)
    \item Entropy remains constant: $\Delta S = 0$ (reversible adiabatic process)
    \item The gas cools because it does work against external pressure without heat input
  \end{itemize}

  \textbf{Key difference}: Isothermal processes maintain temperature via heat exchange; adiabatic 
  processes are thermally isolated, so temperature changes as work is done.

  \item \emph{Sample answer:}\\
  \textbf{Chemical potential}: $\mu$ is the change in Gibbs free energy when one particle is added 
  to the system at constant $T$ and $P$: $\mu = \left(\frac{\partial G}{\partial N}\right)_{T,P}$.

  \textbf{Phase equilibrium condition}: At equilibrium between phases (e.g., liquid and vapor), the 
  chemical potentials must be equal: $\mu_\text{liquid} = \mu_\text{vapor}$.

  \textbf{Physical meaning}:
  \begin{itemize}
    \item If $\mu_\text{liquid} > \mu_\text{vapor}$, particles spontaneously move from liquid to vapor 
    (evaporation) to lower the total free energy
    \item If $\mu_\text{liquid} < \mu_\text{vapor}$, condensation occurs
    \item Equilibrium is reached when $\mu$ is uniform across all phases
  \end{itemize}

  \textbf{Example - Vapor-liquid coexistence}: At a given temperature, the vapor pressure adjusts 
  until $\mu_\text{liquid}(T,P) = \mu_\text{vapor}(T,P)$. This determines the saturation vapor 
  pressure and the coexistence curve in the phase diagram.

\end{enumerate}

\end{document}
